\chapter{Introducción}
En este capítulo se presenta el concepto de oscilador armónico \citep{Biot1851} y su relevancia en el desarrollo de la mecánica cuántica \citep{born1925}.

\section{Antecedentes y contexto}
En el Zhou Li \citep{biot1851tcheou}, obra fundamental de la literatura China,  se documenta el uso de gnomon una vara vertical de aproximadamente 1.6 metros de altura. Mediante la observación de la sombra solar, el gnomon representó la primera estandarización de medición del tiempo y el espacio en Oriente. Este reloj antiguo posee un límite de precisión física de alrededor de 2 minutos, debido al diámetro angular del Sol y el efecto penumbra resultante \citep{kepler1604}. Expresado en terminos de incertidumbre fraccional  \citep{eisenhart1963} esta se estima en el orden de $10^{-2} \text{ a } 10^{-3}$.

De aquella manera rudimentaria de medir el tiempo, casi finalizando el primer cuarto del siglo XXI, nos encontramos ante el reloj atómico Aetherclock OC020 \citep{okaba2024optical}, con una incertidumbre fraccional de $10^{-18}$. Esto equivale a decir que tendría una precisión de 1 segundo cada 30 mil millones de años. Este es el primer reloj de red óptica de estroncio \citep{takamoto2003} comercial lanzado por la empresa \cite{shimadzu2025aetherclock}. Mientras el gnomon usa la sombra solar como un oscilador rudimentario, el Aetherclock implementa los principios del oscilador armónico cuántico que alcanza una precisión de $10^{-18}$, reemplazando la observación  física \citep{Biot1851} por la estabilidad de los niveles de energía atómicos \citep{bohr1913}. 


\section{Motivación y justificación del tema}
El oscilador armónico cuántico (OAC)\citep{Biot1851} constituye, junto con el átomo de hidrógeno, uno de los modelos fundamentales sobre los que se construyó la mecánica cuántica \citep{born1925}. Su origen conceptual se remonta a los primeros intentos por comprender la interacción entre la radiación y la materia a finales del siglo XIX \citep{wien1896, rayleigh1900, jeans1905}, cuando los modelos clásicos resultaron insuficientes para describir correctamente la distribución espectral \citep{wien1896} de la radiación térmica \citep{rayleigh1900}. En este contexto, las investigaciones sobre la distribución espectral evidenciaron las limitaciones de la física clásica, preparando el terreno para la hipótesis de cuantización \citep{planck1901}, quien modeló la materia como un conjunto de resonadores armónicos con energías discretas.


    
La hipótesis de Planck adquirió un significado físico más profundo cuando Einstein la aplicó al estudio del calor específico de los sólidos, interpretando los átomos como osciladores armónicos cuánticos independientes \citep{einstein1907}. Este enfoque permitió explicar la disminución del calor específico a bajas temperaturas, en contraste con las predicciones clásicas basadas en la ley de Dulong y Petit \citep{dulong1819}. Posteriormente, \citet{debye1912} extendió este tratamiento incorporando un espectro continuo de frecuencias vibracionales, consolidando el papel del oscilador armónico en la descripción microscópica de la materia.
	
El desarrollo formal de la mecánica cuántica en la década de 1920 otorgó al oscilador armónico cuántico un rol central en las distintas formulaciones de la teoría. En la mecánica matricial de  \citet{heisenberg1925}, el abandono del concepto de trayectoria clásica y la introducción de magnitudes no conmutativas encuentran en el oscilador armónico uno de sus ejemplos paradigmáticos. Esta formulación fue sistematizada por Born y Jordan (\citeyear{born1925}) y extendida a sistemas generales por Born, Heisenberg y Jordan (\citeyear{bornheisenberg1926}). De manera complementaria, la mecánica ondulatoria de Schrödinger permitió resolver exactamente el oscilador armónico como un problema de valores propios, mostrando que la cuantización de la energía surge de las condiciones matemáticas impuestas a la ecuación de onda \citep{schrodinger1926a, schrodinger1926b}. Finalmente, \citet{dirac1927} demostró que cada modo del campo electromagnético se comporta como un oscilador armónico cuántico, estableciendo este modelo como un elemento fundamental de la teoría cuántica de la radiación .
	
Estas contribuciones consolidaron al oscilador armónico cuántico no solo como un modelo pedagógico, sino como un sistema estructural de la mecánica cuántica, presente en la descripción de una amplia variedad de fenómenos físicos\citep{manuzio_aristoteles_1497}. Su estudio detallado resulta, por tanto, esencial para comprender los fundamentos teóricos de la física cuántica y los principios matemáticos que subyacen a sistemas cuánticos más complejos\citep{dirac1930principles}.


\section{Planteamiento del Trabajo}
El presente Trabajo tiene como objetivo principal analizar el oscilador armónico cuántico como sistema fundamental de la mecánica cuántica, desarrollando de manera explícita su formulación teórica y su resolución matemática. Se busca mostrar cómo este modelo permite introducir conceptos esenciales tales como la cuantización de la energía \citep{schrodinger1926a, schrodinger1926b}, la formulación hamiltoniana \citep{dirac1930}, las relaciones de conmutación \citep{born1925} y el planteamiento del problema cuántico como un problema de valores propios \citep{schrodinger1926a, schrodinger1926b}.
	
%Para ello, el Trabajo parte de un recorrido histórico que contextualiza el surgimiento del oscilador armónico cuántico dentro del desarrollo de la teoría cuántica, apoyándose en los trabajos originales que sentaron sus bases conceptuales. Posteriormente, se aborda el cálculo cuántico del oscilador armónico, que constituye el eje central del documento, poniendo énfasis en la interpretación física de los resultados obtenidos y en su coherencia con los principios generales de la mecánica cuántica.

\section{Estructura del Trabajo}
%El Trabajo se organiza de la siguiente manera. En primer lugar, se presenta un contexto histórico y conceptual del oscilador armónico cuántico, revisando las contribuciones fundamentales que condujeron a su formulación dentro de la mecánica cuántica. A continuación, se desarrolla la descripción clásica del oscilador armónico y su transición hacia el tratamiento cuántico mediante el formalismo hamiltoniano y la ecuación de Schrödinger.
	
%La sección central del trabajo está dedicada al cálculo cuántico del oscilador armónico, donde se obtienen los niveles de energía y las funciones propias del sistema, destacando su significado físico. Finalmente, se discute la relevancia del oscilador armónico cuántico como modelo prototipo y su papel en la comprensión de sistemas cuánticos más complejos.

% REVISAR-----------------------------------------------------------------------------------

A continuación, se presenta la organización del Trabajo:

\begin{itemize}
    \item En el capítulo 2 se presentan los antecedentes que dieron origen a la mecánica cuántica, analizando el problema de la radiación del cuerpo negro y los conceptos de las partículas a campos.
    \item En el capítulo 3 se detallan los objetivos de este Trabajo.
    \item En el capítulo 4 se presenta el cálculo cuántico del oscilador armónico, donde se obtienen los niveles de energía y las funciones propias del sistema, destacando su significado físico.
    \item Finalmente, en el capítulo 5, se discute la relevancia del oscilador armónico cuántico como modelo prototipo y su papel en la comprensión de sistemas cuánticos más complejos.
\end{itemize}

